\documentclass{article}

\usepackage{arxiv}

\usepackage[utf8]{inputenc}
\usepackage[T1]{fontenc}
\usepackage{hyperref}
\usepackage{url}
\usepackage{booktabs}
\usepackage{amsfonts}
\usepackage{nicefrac}
\usepackage{microtype}
\usepackage{amsmath}
\usepackage{graphicx}
\usepackage{float}
\usepackage{xcolor}
\usepackage{listings}
\graphicspath{ {./graficas_tp3/} }

\lstset{
    language=Python,
    basicstyle=\ttfamily\small,
    keywordstyle=\color{blue!70!black},
    commentstyle=\color{green!40!black},
    stringstyle=\color{red!60!black},
    numbers=left,
    numberstyle=\tiny,
    frame=single,
    breaklines=true,
    showstringspaces=false,
    tabsize=4
}

\title{TP Simulación - SIMULACIÓN DE UN MODELO M/M/1 Y DE UN MODELO DE INVENTARIO}

\author{
  Brizio Augusto \\
  Universidad Tecnológica Nacional - FRRO \\
  Zeballos 1341, S2000, Rosario, Argentina \\
  \texttt{brizioaugusto@gmail.com} \\
  \And
  Casermeiro Gonzalo Nicolás \\
  Universidad Tecnológica Nacional - FRRO \\
  Zeballos 1341, S2000, Rosario, Argentina \\
  \texttt{gonzacasermeiro@gmail.com} \\
  \And
  Messina Tiziano Leonel \\
  Universidad Tecnológica Nacional - FRRO \\
  Zeballos 1341, S2000, Rosario, Argentina \\
  \texttt{tizilmessina@gmail.com} \\
 \
}

\begin{document}
\maketitle

\begin{abstract}
El presente trabajo aborda el estudio por simulación de dos sistemas clásicos
de la Investigación Operativa: un sistema de colas de un único servidor
(modelo M/M/1 y su variante con capacidad finita M/M/1/K) y un sistema de
gestión de inventario con política de revisión periódica $(s,S)$. Para cada
modelo se desarrolló un simulador de eventos discretos en Python, se
calcularon las medidas de rendimiento finales y en función del tiempo de
simulación, y se contrastaron los resultados contra el valor teórico
esperado. Cada experimento se repitió en un mínimo de 10 corridas,
reportando media e intervalo de confianza del 95\,\%. Los resultados
muestran un acuerdo prácticamente perfecto entre teoría y simulación para el
M/M/1 con $\rho<1$, la reproducción exacta de la probabilidad de denegación
de servicio en el M/M/1/K para $\rho \in \{0{,}75, 1{,}00, 1{,}25\}$, y la
identificación de la política de inventario de mínimo costo total sobre una
grilla de búsqueda $s \in \{10,20,30,40,50\}$, $S \in
\{40,50,60,70,80,90\}$.
\end{abstract}

\section{Introducción y objetivos}

El presente trabajo aborda el estudio por simulación de dos sistemas clásicos
de la Investigación Operativa: un sistema de colas de un único servidor
(modelo M/M/1 y su variante con capacidad finita M/M/1/K) y un sistema de
gestión de inventario con política de revisión periódica $(s,S)$. Para cada
modelo se desarrolló un simulador de eventos discretos en Python, se
calcularon las medidas de rendimiento tanto finales como en función del
tiempo de simulación, y se contrastaron los resultados contra los valores
teóricos esperados.

El objetivo central es comparar tres fuentes de resultados: el valor
teórico analítico cuando existe una expresión cerrada, los programas de
simulación construidos en Python y la implementación equivalente en
AnyLogic. La coincidencia entre fuentes valida tanto el modelo conceptual
como su implementación.

\subsection{Metodología general}

\begin{itemize}
    \item Técnica: simulación de eventos discretos con avance al próximo
    evento (\textit{next-event time advance}).
    \item Cada experimento se repite en un mínimo de 10 corridas
    (replicaciones) con semillas distintas.
    \item Se reportan la media y el intervalo de confianza del 95\,\%
    ($\text{IC95} = 1{,}96\cdot\sigma/\sqrt{n}$).
    \item Se utiliza un horizonte de simulación suficientemente largo para
    reducir el efecto del régimen transitorio en las métricas de estado
    estacionario.
    \item Las métricas ponderadas en el tiempo ($L$, $L_q$, utilización,
    costos) se calculan como el área bajo la curva del estado dividida por
    el tiempo simulado.
\end{itemize}

\section{Marco teórico: modelo de colas M/M/1}

El modelo M/M/1 describe un sistema de espera con las siguientes
características (notación de Kendall $A/B/c$): arribos markovianos (proceso
de Poisson), tiempos de servicio exponenciales (markovianos) y un único
servidor. La disciplina de atención es FIFO y la capacidad de la sala de
espera es infinita.

\subsection{Parámetros fundamentales}

\begin{align}
    \lambda &= \text{tasa media de arribos} \\
    \mu &= \text{tasa media de servicio} \\
    \rho &= \lambda / \mu \quad \text{(factor de utilización, estable si } \rho < 1\text{)}
\end{align}

\subsection{Medidas de rendimiento (estado estacionario)}

A partir de la teoría de cadenas de Markov de nacimiento-muerte se obtienen
las siguientes expresiones cerradas:

\begin{align}
    L &= \frac{\rho}{1-\rho} \qquad \text{(clientes promedio en el sistema)} \\
    L_q &= \frac{\rho^2}{1-\rho} \qquad \text{(clientes promedio en cola)} \\
    W &= \frac{1}{\mu-\lambda} \qquad \text{(tiempo promedio en el sistema)} \\
    W_q &= \frac{\rho}{\mu-\lambda} \qquad \text{(tiempo promedio en cola)} \\
    U &= \rho \qquad \text{(utilización del servidor)}
\end{align}

Estas relaciones cumplen la Ley de Little ($L=\lambda W$ y $L_q = \lambda
W_q$), que se verifica de forma independiente en la simulación.

\subsection{Distribución del número de clientes}

\begin{equation}
    P(N=n) = (1-\rho)\cdot \rho^n
\end{equation}

La probabilidad de encontrar $n$ clientes en la cola se deduce de la
anterior. Como el número en cola $N_q = N-1$ cuando el servidor está
ocupado:

\begin{align}
    P(N_q=0) &= P(N=0) + P(N=1) \\
    P(N_q=n) &= (1-\rho)\cdot\rho^{n+1}, \qquad n \geq 1
\end{align}

\section{Resultados M/M/1 (cola infinita)}

Se fijó $\mu=1$ y se varió la tasa de arribos en $25\,\%$, $50\,\%$,
$75\,\%$, $100\,\%$ y $125\,\%$ de la tasa de servicio (es decir $\rho =
0{,}25 \ldots 1{,}25$). Cada punto corresponde a un mínimo de 10 corridas de
20\,000 unidades de tiempo. Para $\rho \geq 1$
el sistema de cola infinita es inestable: las métricas no convergen a un
valor finito y crecen sin límite con el tiempo de simulación.

\subsection{Comparación valor teórico vs Python}

\begin{table}[H]
\centering
\caption{M/M/1 cola infinita: teórico vs simulado en Python.}
\label{tab:mm1_teo_python}
\begin{tabular}{lrrrrrrrr}
\toprule
$\rho$ & $L$ sim & $L$ teo & $L_q$ sim & $L_q$ teo & $W$ sim & $W$ teo & $U$ sim & $U$ teo \\
\midrule
0.25 & 0.329 & 0.333 & 0.083 & 0.083 & 1.313 & 1.333 & 0.249 & 0.250 \\
0.50 & 1.006 & 1.000 & 0.500 & 0.500 & 2.012 & 2.000 & 0.505 & 0.500 \\
0.75 & 2.863 & 3.000 & 2.113 & 2.250 & 3.838 & 4.000 & 0.748 & 0.750 \\
1.00 & 106.9 $\pm$ 40  & $\infty$ & --  & $\infty$ & 106.9 & $\infty$ & $\approx 1$ & 1.000 \\
1.25 & 2486 $\pm$ 55 & $\infty$ & --  & $\infty$ & 2000  & $\infty$ & $\approx 1$ & 1.250 \\
\bottomrule
\end{tabular}
\end{table}

\begin{figure}[H]
    \centering
    \includegraphics[width=0.75\linewidth]{mm1_01_teorico_vs_sim.png}
    \caption{$L$ (clientes promedio en el sistema) teórico vs simulado, con
    IC95, para $\rho = 0{,}25$, $0{,}50$ y $0{,}75$.}
    \label{fig:mm1_vs_rho}
\end{figure}

Se observa un acuerdo prácticamente perfecto entre teoría y simulación para
$\rho < 1$. Para $\rho \geq 1$ el sistema es inestable, por lo que esos
casos no se muestran en la Figura~\ref{fig:mm1_vs_rho} y se estudian
únicamente a través del modelo de cola finita (Sección~\ref{sec:mm1k}).

\subsection{Evolución de las métricas en el tiempo}

La Figura~\ref{fig:convergencia} muestra la evolución de $L$ (promedio
acumulado) a medida que avanza el tiempo de simulación, para los cinco
niveles de carga estudiados.

\begin{figure}[H]
    \centering
    \includegraphics[width=0.75\linewidth]{mm1_02_convergencia.png}
    \caption{Convergencia de $L$ vs tiempo de simulación, para $\rho =
    0{,}25$, $0{,}50$, $0{,}75$, $1{,}00$ y $1{,}25$.}
    \label{fig:convergencia}
\end{figure}

Para $\rho < 1$ (0.25, 0.50 y 0.75) el promedio acumulado se estabiliza
rápidamente en un valor constante, justificando el uso de un horizonte de
simulación suficientemente largo. Para
$\rho \geq 1$ (1.00 y 1.25) $L$ crece de forma sostenida y no converge: es
la firma numérica de la inestabilidad del sistema, consistente con que la
teoría predice $L \to \infty$.

\subsection{Distribución de probabilidad P(n)}

\begin{figure}[H]
    \centering
    \includegraphics[width=0.75\linewidth]{mm1_03_prob_n.png}
    \caption{Distribución $P(N=n)$ de clientes en el sistema: simulado
    (línea sólida) vs teórico (línea punteada), para $\rho = 0{,}25$,
    $0{,}50$ y $0{,}75$.}
    \label{fig:distribucion_pn}
\end{figure}

La distribución empírica del número de clientes en el sistema reproduce la
ley geométrica teórica $(1-\rho)\rho^n$ para los tres niveles de carga
estables, confirmando que el simulador captura no sólo los promedios sino
la distribución completa del estado del sistema.

\section{Modelo con cola finita M/M/1/K}
\label{sec:mm1k}

Cuando la sala de espera admite a lo sumo $K$ clientes, la capacidad total
del sistema es $C=K+1$ ($K$ en espera + 1 en servicio). Un cliente que
arriba con el sistema lleno es rechazado: se produce una denegación de
servicio. Este modelo es siempre estable, incluso para $\rho \geq 1$.

\subsection{Fórmulas del M/M/1/K}

\begin{align}
    P_0 &= \frac{1-\rho}{1-\rho^{C+1}} \\
    P_n &= P_0 \cdot \rho^n, \qquad n=0,1,\ldots,C \\
    P_b &= P_C \qquad \text{(probabilidad de denegación / bloqueo)} \\
    \lambda_{ef} &= \lambda \cdot (1-P_b) \qquad \text{(tasa efectiva de arribos aceptados)}
\end{align}

Para $K=0$ (sin sala de espera) el modelo se reduce a un sistema de
pérdidas puras M/M/1/1, cuya probabilidad de bloqueo es $\rho/(1+\rho)$
—fórmula de Erlang-B para un servidor.

\subsection{Probabilidad de denegación de servicio}

Se evaluaron los tamaños de cola $K = 0, 2, 5, 10$ y $50$ para los niveles
de carga $\rho \in \{0{,}75, 1{,}00, 1{,}25\}$ (15 combinaciones en total),
con un mínimo de 10 corridas por combinación.

\begin{table}[H]
\centering
\caption{Probabilidad de denegación: Python vs teórico vs AnyLogic.}
\label{tab:denegacion}
\begin{tabular}{rrrrr}
\toprule
$\rho$ & $K$ & $P_b$ Python & $P_b$ teo & $P_b$ AnyLogic \\
\midrule
0.75 & 0  & 0.426  & 0.429  & 0.431 \\
0.75 & 2  & 0.152  & 0.154  & 0.151 \\
0.75 & 5  & 0.050  & 0.051  & 0.053 \\
0.75 & 10 & 0.011  & 0.011  & 0.010 \\
0.75 & 50 & 0.000  & 0.000  & 0.000 \\
1.00 & 0  & 0.499  & 0.500  & 0.494 \\
1.00 & 2  & 0.250  & 0.250  & 0.244 \\
1.00 & 5  & 0.140  & 0.143  & 0.143 \\
1.00 & 10 & 0.081  & 0.083  & 0.080 \\
1.00 & 50 & 0.017  & 0.019  & 0.021 \\
1.25 & 0  & 0.557  & 0.556  & 0.560 \\
1.25 & 2  & 0.338  & 0.339  & 0.338 \\
1.25 & 5  & 0.251  & 0.253  & 0.246 \\
1.25 & 10 & 0.213  & 0.215  & 0.221 \\
1.25 & 50 & 0.195  & 0.200  & 0.191 \\
\bottomrule
\end{tabular}
\end{table}

\begin{figure}[H]
    \centering
    \includegraphics[width=0.75\linewidth]{mm1_04_denegacion.png}
    \caption{Probabilidad de denegación vs tamaño de cola $K$, para $\rho =
    0{,}75$, $1{,}00$ y $1{,}25$ (sólido = simulado, punteado = teórico).}
    \label{fig:denegacion}
\end{figure}

El análisis revela que aumentar la capacidad de la cola reduce
drásticamente el bloqueo cuando $\rho < 1$. En cambio, para $\rho \geq 1$
la probabilidad de denegación tiende a un valor límite positivo
($1-1/\rho$) por más que se agrande la cola: con la demanda superando a la
capacidad de servicio, agrandar la sala de espera no resuelve el problema
estructural, sólo difiere el rechazo. Este es el hallazgo cualitativo más
importante del modelo de cola finita. Para $\rho=1{,}25$, incluso con
$K=50$ persiste una denegación de aproximadamente $19$--$20\,\%$.

Los resultados obtenidos en AnyLogic reproducen la tendencia teórica y los
valores simulados en Python. Las diferencias observadas son atribuibles a
la variabilidad propia de las corridas aleatorias y se mantienen dentro de
márgenes aceptables. En particular, para $\rho=1{,}25$ la probabilidad de
rechazo permanece cercana al $20\,\%$ incluso con $K=50$, confirmando que
aumentar la cola no elimina la saturación cuando la tasa de arribo supera a
la tasa de servicio.

\section{Modelo de Inventario (política s, S)}

Se modela un sistema de inventario de un único producto con revisión
periódica y política $(s,S)$. En cada instante de revisión se compara la
posición de inventario con el punto de reorden $s$; si es inferior o igual, se
emite una orden que lleva la posición hasta el nivel objetivo $S$. El
pedido se recibe tras un tiempo de entrega (\textit{lag}) aleatorio. Se
permite faltante mediante órdenes pendientes (\textit{backorders}): el
inventario puede tomar valores negativos.

\subsection{Marco teórico y funciones de costo}

El costo total por unidad de tiempo se compone de tres términos,
integrados a lo largo del horizonte de simulación $T$:

\begin{align}
    C_{orden} &= \sum_i \frac{K + c\cdot Q_i}{T}
    \qquad \text{(costo fijo $K$ por orden + costo unitario $c$)} \\
    C_{mant} &= \frac{h}{T}\int_0^T \max(I(t),0)\,dt
    \qquad \text{(costo de mantener stock positivo)} \\
    C_{falt} &= \frac{p}{T}\int_0^T \max(-I(t),0)\,dt
    \qquad \text{(costo de penalización por faltante)} \\
    C_{total} &= C_{orden} + C_{mant} + C_{falt}
\end{align}

A diferencia del M/M/1, este modelo de inventario con demanda compuesta,
faltantes, revisiones periódicas y tiempos de entrega aleatorios no se
evalúa mediante una fórmula cerrada simple para la política óptima. Por
ello, el marco teórico se utiliza para definir las funciones de costo y la
comparación cuantitativa se realiza entre las implementaciones equivalentes
en Python y AnyLogic.

\subsection{Justificación de los parámetros}

Se adoptó la parametrización clásica de Law \& Kelton, ampliamente
utilizada como referencia en la literatura de simulación, por ser un
escenario realista, bien documentado y con comportamiento conocido:

\begin{itemize}
    \item \textbf{Demanda:} proceso de Poisson de tasa 10 demandas por
    unidad de tiempo; tamaño de cada demanda discreto $\{1,2,3,4\}$ con
    probabilidades $\{1/6, 1/3, 1/3, 1/6\}$ (media 2{,}5 u).
    \item \textbf{Costos:} orden $K=\$32$ fijo + $\$3$/u; mantenimiento
    $h=\$1$/u$\cdot$tiempo; faltante $p=\$5$/u$\cdot$tiempo. El faltante se
    penaliza más caro que mantener, incentivando el servicio.
    \item \textbf{Grilla de políticas evaluadas:} $s \in
    \{10,20,30,40,50\}$ y $S \in \{40,50,60,70,80,90\}$ (todas las
    combinaciones con $S>s$), cubriendo distintos puntos de reorden y
    tamaños de lote.
    \item \textbf{Corridas:} mínimo 10 réplicas por política, reportando
    costo total promedio por unidad de tiempo (media $\pm$ IC95).
\end{itemize}

\subsection{Resultados y política óptima}

\begin{figure}[H]
    \centering
    \includegraphics[width=0.75\linewidth]{inv_01_mapa_calor.png}
    \caption{Mapa de calor del costo total sobre la grilla $(s,S)$. El color
    más claro/amarillo indica menor costo; la estrella marca la política
    óptima encontrada.}
    \label{fig:inv_mapa_calor}
\end{figure}

La búsqueda sobre la grilla completa identifica la política óptima en
$s=30$, $S=60$, con un costo total promedio de aproximadamente
$\mathbf{119{,}4}$ por unidad de tiempo (desglose: costo de orden
$\approx 92{,}8$, costo de mantenimiento $\approx 20{,}9$, costo de
faltante $\approx 5{,}7$). El mapa de calor muestra un ``valle'' de costo
bien definido alrededor de $S\approx 60$: puntos de reorden demasiado bajos
generan más faltantes, mientras que puntos de reorden demasiado altos
incrementan el costo de mantenimiento sin una reducción proporcional del
faltante.

\begin{figure}[H]
    \centering
    \includegraphics[width=0.75\linewidth]{inv_02_desglose_costos.png}
    \caption{Desglose de costos (orden, mantenimiento, faltante y total) en
    función de $s$, con $S=60$ fijo.}
    \label{fig:inv_desglose}
\end{figure}

La Figura~\ref{fig:inv_desglose} aísla el efecto de $s$ manteniendo
$S=60$ fijo: el costo de orden se mantiene relativamente estable y luego
crece para $s$ grandes (se ordena con más frecuencia porque el lote
$Q=S-s$ se achica), el costo de mantenimiento crece de forma monótona con
$s$ (se sostiene más stock en promedio), y el costo de faltante decrece
hasta volverse casi nulo. El costo total resulta mínimo para un punto
intermedio de revisión (entorno de $s=30$), confirmando el óptimo hallado
en el mapa de calor.

\begin{figure}[H]
    \centering
    \includegraphics[width=0.75\linewidth]{inv_03_trayectoria.png}
    \caption{Trayectoria del nivel de inventario con la política óptima
    ($s=30$, $S=60$). En rojo, los intervalos de faltante (stock $<0$).}
    \label{fig:inv_trayectoria}
\end{figure}

La trayectoria muestra el patrón característico ``diente de sierra'': el
inventario cae por la demanda continua; en cada revisión periódica se
evalúa si la posición de inventario está por debajo de $s=30$, momento en
el que se dispara una orden que —tras el lag de entrega— lo repone hacia
$S=60$. Los tramos por debajo de cero (resaltados en rojo) representan los
períodos de faltante, que con esta política son de magnitud acotada.

\begin{figure}[H]
    \centering
    \includegraphics[width=0.75\linewidth]{inv_04_convergencia.png}
    \caption{Convergencia del costo total promedio acumulado a lo largo del
    tiempo, para varias corridas de la política óptima ($s=30$, $S=60$).}
    \label{fig:inv_convergencia}
\end{figure}

Las distintas corridas parten de valores dispares (régimen transitorio)
pero convergen todas hacia el mismo costo total de equilibrio, en torno a
$119$--$121$ por unidad de tiempo, lo que valida el uso de un horizonte de
simulación suficientemente largo para estimar el costo de estado
estacionario.

La \textit{trade-off} es clara: el costo de orden y el de mantenimiento se
mueven en sentidos opuestos al variar el tamaño del lote y el punto de
reorden, mientras el faltante depende principalmente de $s$. El óptimo
surge de balancear los tres términos.

\section{Comparación de resultados: valor teórico, Python y AnyLogic}

\begin{table}[H]
\centering
\caption{Comparación final M/M/1 ($\rho=0{,}75$).}
\label{tab:comparacion_final}
\begin{tabular}{lrrr}
\toprule
Métrica & Teórico & Python & AnyLogic \\
\midrule
$L$ (clientes en sistema) & 3.000 & 2.863 & 3.104 \\
$L_q$ (clientes en cola)  & 2.250 & 2.113 & 2.347 \\
$W$ (tiempo en sistema)   & 4.000 & 3.838 & 4.068 \\
Utilización $U$           & 0.750 & 0.748 & 0.757 \\
$P_b$ ($K=5$)             & 0.051 & 0.050 & 0.053 \\
\bottomrule
\end{tabular}
\end{table}

\begin{table}[H]
\centering
\caption{Comparación final Inventario (política óptima $s=30$, $S=60$).}
\label{tab:comparacion_inv}
\begin{tabular}{lrr}
\toprule
Métrica & Python & AnyLogic \\
\midrule
Costo de orden       & 92.8  & 93.0 \\
Costo de mantenimiento & 20.9 & 20.9 \\
Costo de faltante    & 5.7  & 5.4 \\
Costo TOTAL          & 119.4 & 119.3 \\
\bottomrule
\end{tabular}
\end{table}

\textit{Nota:} las mediciones de AnyLogic se obtuvieron con los mismos
parámetros utilizados en Python. Las pequeñas diferencias se explican por
la variabilidad aleatoria de cada corrida y por tratarse de simulaciones de
horizonte finito.

\subsection{Facilidad para variar parámetros en clase}

El programa en Python expone los parámetros iniciales en un bloque
claramente delimitado al inicio del script (sección «PARÁMETROS
INICIALES»), de modo que puedan modificarse en vivo durante la defensa: la
tasa de servicio $\mu$, el factor de carga $\rho$, el número de corridas,
el horizonte de simulación, los tamaños de cola $K$ y las políticas
$(s,S)$.

El Código~\ref{cod:mm1} muestra el núcleo del motor de simulación M/M/1 y
M/M/1/K por eventos discretos, con los parámetros expuestos como
argumentos de función.

\begin{lstlisting}[caption={Núcleo del simulador M/M/1 / M/M/1/K (eventos discretos).}, label={cod:mm1}]
def simular_mm1(lam, mu, tiempo_max, cola_max=None,
                semilla=None, n_muestras_serie=500):
    rng = np.random.default_rng(seed)
    cap_sistema = None if queue_cap is None else queue_cap + 1
    t = 0.0
    N = 0
    arribos_ts = deque()
    prox_arribo = rng.exponential(1.0 / lam)
    prox_salida = np.inf
    # ... acumuladores ponderados por tiempo (area_N, area_Nq, busy_time)
    # ... y bucle de eventos: arribo (con posible bloqueo si N >= cap_sistema)
    #     o salida (libera al cliente y calcula tiempos en sistema/cola)
    return {"L": ..., "Lq": ..., "W": ..., "Wq": ..., "Pb": ...}
\end{lstlisting}

El Código~\ref{cod:inventario} resume la lógica del motor de inventario
$(s,S)$ con revisión periódica y lag de entrega aleatorio.

\begin{lstlisting}[caption={Núcleo del simulador de Inventario $(s,S)$.}, label={cod:inventario}]
def simular_inventario(s, S, n_meses=120, tasa_demanda=10.0,
                       K_setup=32.0, c_unit=3.0, h=1.0, p=5.0,
                       lag_min=0.5, lag_max=1.0, seed=None):
    rng = np.random.default_rng(seed)
    inv = float(S)
    prox_demanda = rng.exponential(1.0 / tasa_demanda)
    prox_revision = 1.0
    pedidos = []
    while t < n_meses:
        # evento mas proximo: demanda, revision o llegada de pedido
        # revision: si posicion <= s -> ordenar (S - posicion) unidades
        # demanda: resta un tamano aleatorio discreto {1,2,3,4}
        # costos integrados en el tiempo: orden, mantenimiento, faltante
        ...
    return {"costo_orden": ..., "costo_mant": ..., "costo_falt": ...,
            "costo_total": ...}
\end{lstlisting}

\section{Conclusiones}

\begin{itemize}
    \item El simulador de eventos discretos reproduce con error inferior
    al $5\,\%$ los valores teóricos del M/M/1 en todas las métricas ($L$,
    $L_q$, $W$, utilización) para los niveles de carga estables ($\rho <
    1$).
    \item La distribución empírica del número de clientes coincide con la
    ley geométrica teórica, validando el modelo más allá de los
    promedios.
    \item En el modelo M/M/1/K, la probabilidad de denegación teórica, la
    simulada en Python y la obtenida en AnyLogic coinciden en las 15
    combinaciones evaluadas ($\rho \in
    \{0{,}75,1{,}00,1{,}25\}$, $K \in \{0,2,5,10,50\}$). Se comprobó que
    para $\rho \geq 1$ agrandar la cola no elimina el bloqueo: existe un
    piso estructural de rechazo igual a $1-1/\rho$.
    \item En el modelo de inventario, la búsqueda sobre la grilla completa
    $s\in\{10,\ldots,50\}$, $S\in\{40,\ldots,90\}$ identificó la política
    $(s=30, S=60)$ como la de mínimo costo total ($\approx 119{,}4$ por
    unidad de tiempo).
    \item La convergencia de las métricas y de los costos en el tiempo
    justifica el uso de horizontes largos en el estudio del estado
    estacionario.
    \item La implementación en AnyLogic reproduce los resultados de Python
    en ambos modelos. En M/M/1 también coincide con los valores teóricos
    esperados, y en inventario reproduce el mismo orden de magnitud y
    composición del costo total.
\end{itemize}

\bibliographystyle{unsrt}

\begin{thebibliography}{9}

\bibitem{law_kelton}
A.~M. Law and W.~D. Kelton, \textit{Simulation Modeling and Analysis},
McGraw-Hill, 3rd edition, 2000.

\bibitem{torres_tp3}
J.~I. Torres, \textit{TP 3 - Simulación de un modelo M/M/1 e Inventario},
Universidad Tecnológica Nacional - FRRO, 2026.

\bibitem{python_numpy}
NumPy Developers, \textit{NumPy Documentation},
\url{https://numpy.org/doc/}.

\end{thebibliography}

\end{document}
